% Created 2010-04-19 Mon 08:15
\documentclass[11pt]{article}
\usepackage[utf8]{inputenc}
\usepackage[T1]{fontenc}
\usepackage{fixltx2e}
\usepackage{graphicx}
\usepackage{longtable}
\usepackage{float}
\usepackage{wrapfig}
\usepackage{soul}
\usepackage{t1enc}
\usepackage{textcomp}
\usepackage{marvosym}
\usepackage{wasysym}
\usepackage{latexsym}
\usepackage{amssymb}
\usepackage{hyperref}
\tolerance=1000
\providecommand{\alert}[1]{\textbf{#1}}

\title{software overview}
\author{Giulio Bottazzi}
\date{19 April 2010}

\begin{document}

\maketitle

\setcounter{tocdepth}{3}
\tableofcontents
\vspace*{1cm}

Below you can find part of the software I wrote for my research and
teaching activity. It is a quite heterogeneous collection of programs,
both in terms of scope and quality.

All these programs have been developed, tested and used in Linux
(different flavors). They should compile and run, with none or minor
modifications, on any Unix platform. The porting to different OSes,
however, can require major rewriting and has never been tested.

Please notice that all the software described in these pages has been
mainly written for personal use, and is distributed under the \href{http://www.gnu.org/copyleft/gpl.html}{General Public License} in the hope it could be of help to other people. It is
provided ``as is'' without express or implied warranty to the extent
permitted by applicable law.

\section{Alka-linux}
\label{sec-1}


Alka-linux is a linux distribution aimed at providing a simple and
user friendly environment for the econometric and economic
analysis. The OS can be booted from a CD and don't require any
previous installation. Light-wight software allow the system to run
smoothly also on older machines. See the project \href{http://www.alka-linux.sssup.it/}{web page}.
\section{gbutils}
\label{sec-2}


A set of statistical utilities. They are intended as an help to people
using gnuplot, indeed their output is designed in a format suitable to
be sent to gnuplot for plotting. They take ASCII data files as
input. Utilities to obtain binned density (histogram), kernel
estimated density, cumulative distribution and binned statistics from
data are included. All the programs are written in C.  See the \href{http://www.cafed.sssup.it/software/gbutils/gbutils.html}{web page}.
\section{subbotools}
\label{sec-3}


A set of programs written to deal with the Subbotin and antisymmetric
Subbotin distributions. Programs for maximum likelihood estimation of
parameters are included, together with generators of random
numbers. All the programs are written in C.  See the \href{http://www.cafed.sssup.it/software/subbotools/index.html}{web page}.
\section{YAFiMM}
\label{sec-4}


This package contains programs providing a simulation framework for
the analysis of a specific financial market model. For the description
of the model, see the working paper \href{http://www.sssup.it/%7ELEM/WPLem/2002-10.html}{A Simple Micro-Model of Market Dynamics Part I: The ``Homogenous Agents'' Deterministic Limit} in the WP
page. All the programs are written in C++. See the \href{http://www.cafed.sssup.it/software/yafimm/index.html"}{web page}.
\section{bose}
\label{sec-5}


This package contains various programs related to the Bottazzi-Secchi
model of industry dynamics described in \href{http://www.sssup.it/%7ELEM/WPLem/2002-20.html}{On the Laplace Distribution of Firms Growth Rates}. Some of the programs are written in Python, some
in C. See the \href{http://www.cafed.sssup.it/software/bose/index.html"}{web page.}
\section{NAIF}
\label{sec-6}


NAIF (Numerical Analysis of Investment Functions) is a program to for
simulating the trading activities of an heterogeneous group og agents
in a simple two-good economy. See the \href{http://www.cafed.sssup.it/software/naif/naif.html"}{web page}.
\section{bdr}
\label{sec-7}


This package contains the file needed to perform simulations of the
model of industrial dynamics described in G.Bottazzi, G.Dosi and
G.Rocchetti Modes of Knowledge Accumulation, Entry Regimes and
Patterns of Industrial Evolution Industrial and Corporate Change 10,3
(2001). All the programs are written in Python and require various
packages and modules that are freely and publicly available.
\section{multimin}
\label{sec-8}


  \href{file://.software/multimin/miltimin.html}{multimin} is an interface to the \href{http://www.gnu.org/software/gsl/}{GNU Scientific Library} minimization
  routines. It is a simple C function that accepts the function to be
  minimized, possibly its derivatives, the initial conditions and the
  minimization options as formal parameters.
\section{spanning tree}
\label{sec-9}


  \href{file://.software/spanntree/spanntree.tar.gx}{spanntree} is a collection of routines for the calculation of maximum
  and minimum spanning tree of weighted graphs. It can be used for
  directed or undirected graph. In case the graph is made of several
  disjoint pieces, the routines return a spanning forest (i.e. the
  union of several spanning tree).

\end{document}